\documentclass[11pt]{amsart}
%%%%%%%%%%%%%%%%%%%%%%%%%%%%%% PACKAGES %%%%%%%%%%%%%%%%%%%%%%%%%%%%%%
\usepackage{amsfonts,amsmath,latexsym,amssymb,verbatim,amsbsy,amsthm}
\usepackage{url}
\usepackage[shortlabels]{enumitem}
\usepackage{caption}
\usepackage{float}
\usepackage{wrapfig}
\usepackage{graphicx}
\usepackage{bigints}
\usepackage{harpoon}
\usepackage{physics}
\usepackage{pythonhighlight}
\usepackage[top=.5in, bottom=.5in, left=.69in, right=.69in]{geometry}
%%%%%%          <<== TEMPLATE
%%%%%%%%%%%%%%%%%%%%%%%%%%%%%% SHORTCUTS %%%%%%%%%%%%%%%%%%%%%%%%%%%%%%
\def \lp {\left(}
\def \rp {\right)}
%%%%%%%%%%%%%%%%%% GREEKS %%%%%%%%%%%%%%%%%%
%%%%%% LOWER %%%%%%
\def \a {\alpha}
\def \b {\beta}
\def \g {\gamma}
\def \d {\delta}
\def \e {\varepsilon}
\def \f {\phi}
\def \ff {\varphi}
\def \l {\lambda}
\def \n {\nabla}
\def \s {\sigma}
\def \t {\tau}
\def \th {\theta}
\def \o {\omega}
\def \w {\omega}
\def \m {\mu}
\def \p {\rho}
\def \ps {\psi}
\def \z {\zeta}
\def \ee {\epsilon}
%%%%%% UPPER %%%%%%
\def \D {\Delta}
\def \G {\Gamma}
\def \F {\Phi}
\def \L {\Lambda}
\def \Th {\Theta}
\def \O {\Omega}
\def \Ps {\Psi}
\def \R {\mathcal{R}}

%%%%%%%%%%%%%%%%%% MATH %%%%%%%%%%%%%%%%%%
%%%%%% COMMANDS %%%%%%
\newcommand{\pd}[2]{\frac{\partial{#1}}{\partial{#2}}}
\newcommand{\ppd}[2]{\frac{\partial^2{#1}}{\partial{#2}^2}}
\newcommand{\drd}[2]{\frac{\text{d}{#1}}{\text{d}{#2}}}
\newcommand{\DrD}[2]{\frac{\text{D}{#1}}{\text{D}{#2}}}
\newcommand{\ans}[1]{\text{\fbox{$#1$}}}
\newcommand{\Rat}[1]{\mathcal{R}_{r,{#1}}}
\newcommand{\Eq}[1]{\begin{equation}{#1}\end{equation}}
\newcommand{\ord}[1]{\mathcal{O}\left({#1}\right)}
%%%%%% SIMPLE %%%%%%
\def \Ra {\Rightarrow}
\def \ra {\rightarrow}
\def \q {\partial}
\def \rb {\rangle}
\def \lb {\langle}
\def \tf {\therefore}
\def \bc {\because}
\def \sin {\text{sin}}
\def \cos {\text{cos}}
\def \tan {\text{tan}}
\def \cot {\text{cot}}
\def \sep {,\qquad}


\makeatletter
\def\widebreve{\mathpalette\wide@breve}
\def\wide@breve#1#2{\sbox\z@{$#1#2$}%
     \mathop{\vbox{\m@th\ialign{##\crcr
\kern0.08em\brevefill#1{0.8\wd\z@}\crcr\noalign{\nointerlineskip}%
                    $\hss#1#2\hss$\crcr}}}\limits}
\def\brevefill#1#2{$\m@th\sbox\tw@{$#1($}%
  \hss\resizebox{#2}{\wd\tw@}{\rotatebox[origin=c]{90}{\upshape(}}\hss$}
\makeatletter

\def \ol {\overline}
\def \til {\tilde}
\def \wt {\widetilde}
\def \bre {\breve}
\def \wb {\widebreve}
\def \wh {\widehat}

%%%%%%%%%%%%%%%%%% UNITS %%%%%%%%%%%%%%%%%%
\def \ms {\frac{\text{m}}{\text{s}}}
\def \acc {\frac{\text{m}}{\text{s}^2}}
\def \den {\frac{\text{kg}}{\text{m}^3}}
\def \pres {\frac{\text{kg}}{\text{m}^2}}
\def \nm {\frac{\text{N}}{\text{m}}}
\def \tr {\text{Tr}}
\def \gas {\frac{\text{J}}{\text{kg}\cdot\text{K}}}
\def \kel {\text{K}}
\def \cel {^o\text{C}}
\def \kpa {\text{kPa}}
\def \atm {\text{atm}}
\def \mf {\frac{\text{kg}}{\text{s}}}
\def \deg {^\text{o}}
\def \rad {\text{rad}}
\def \Re {\text{Re}}
\def \Pr {\text{Pr}}
\def \Nu {\text{Nu}}
\def \con {\frac{\text{W}}{\text{m}\cdot\text{K}}}

\usepackage{mhchem}

%%%%%%%%%%%%%%%%%%%%%%%%%%%%%%%%%%%% DOCUMENT %%%%%%%%%%%%%%%%%%%%%%%%%%%%%%%%%%%%
\begin{document}
\title{2025/10/29 Chamorro Proposal Stuff}
\author{The Navier-Stoked}
\maketitle

%%%%%%%%%%%%%%%%%%%%%%%%%%%%%%%%%%%%%%%%%%%%%%%%%%%%%%%%%%%%%%%%%%%%%%%%%%%%%%%%
\section{Concepts for Wall Actuation in Boundary-Layer Flows}
%%%%%%%%%%%%%%%%%%%%%%%%%%%%%%%%%%%%%%%%%%%%%%%%%%%%%%%%%%%%%%%%%%%%%%%%%%%%%%%%

Two experimental actuation strategies are proposed for the study of boundary-layer response to surface motion. The first introduces controlled wall--normal deformation of the test surface, directly altering the near-wall shear distribution. The second applies spanwise contraction or expansion of the lateral boundaries to impose cross-flow strain and modify the three-dimensional structure of turbulence. Together, these approaches provide a framework to assess directional sensitivity of wall-bounded turbulence to imposed wall motion.

%%%%%%%%%%%%%%%%%%%%%%%%%%%%%%%%%%%%%%%%%%%%%%%%%%%%%%%%%%%%%%%%%%%%%%%%%%%%%%%%
\subsection{Concept~A: Wall--Normal Surface Actuation}
%%%%%%%%%%%%%%%%%%%%%%%%%%%%%%%%%%%%%%%%%%%%%%%%%%%%%%%%%%%%%%%%%%%%%%%%%%%%%%%%

In this concept, the solid boundary is rendered dynamically deformable by an embedded array of actuators producing time-dependent displacements in the wall-normal direction. By controlling these actuators, both standing and traveling deformation waves can be generated, enabling targeted excitation of specific turbulent scales.


The imposed deformation modifies the instantaneous wall-normal velocity boundary condition and locally perturbs the pressure field. 

\begin{itemize}
    \item Quantify the effect of wall-normal surface motion on Reynolds stress distribution and turbulent kinetic energy near the wall.
    \item Examine phase relationships between the actuation waveform and the natural bursting cycle.
    \item Identify control regimes where deformation waves promote energy redistribution across turbulent scales.
\end{itemize}

%%%%%%%%%%%%%%%%%%%%%%%%%%%%%%%%%%%%%%%%%%%%%%%%%%%%%%%%%%%%%%%%%%%%%%%%%%%%%%%%
\subsection{Concept~B: Spanwise Contraction and Lateral Actuation}
%%%%%%%%%%%%%%%%%%%%%%%%%%%%%%%%%%%%%%%%%%%%%%%%%%%%%%%%%%%%%%%%%%%%%%%%%%%%%%%%

The second approach investigates the influence of time-varying lateral confinement on boundary-layer structure.
Here, the spanwise boundaries are driven with synchronized harmonic motion that periodically contracts and expands the effective channel width. Lateral actuation introduces a controlled, time-dependent spanwise strain field. This strain alters both the mean pressure distribution and the instantaneous orientation of large-scale motions within the boundary layer. 
\begin{itemize}
    \item Characterize how spanwise strain affects the organization and strength of streamwise vortices and low-speed streaks.
    \item Compare lateral and wall-normal forcing to determine anisotropic response of turbulence to multi-axis wall motion.
    \item Explore potential mechanisms for turbulent drag reduction or redistribution of Reynolds stresses through dynamic confinement.
\end{itemize}

%%%%%%%%%%%%%%%%%%%%%%%%%%%%%%%%%%%%%%%%%%%%%%%%%%%%%%%%%%%%%%%%%%%%%%%%%%%%%%%%
\subsection{Complementary Scope}
%%%%%%%%%%%%%%%%%%%%%%%%%%%%%%%%%%%%%%%%%%%%%%%%%%%%%%%%%%%%%%%%%%%%%%%%%%%%%%%%

The two actuation concepts are intentionally orthogonal in their mode of forcing: 
Concept~A modifies the local wall-normal boundary condition and directly exchanges momentum with the near-wall fluid, 
while Concept~B

\noindent\rule{\textwidth}{0.4pt}

5cmx60cm 

should be fine with both ``orientations of the wind tunnel'' --- approximate boundary layer thickness allows for valid testing






\end{document}